
\section{Un diagrama de equilibrio no es una historia de proceso}
Un diagrama de fases representa equilibrios para composición y temperatura bajo las condiciones del modelo. Una línea de enlace permite identificar composiciones de fases coexistentes; la regla de la palanca proporciona sus fracciones en masa cuando el eje de composición usa esa base. Para una aleación global $C_0$ entre $C_\alpha$ y $C_\beta$,
\begin{equation}w_\alpha=\frac{C_\beta-C_0}{C_\beta-C_\alpha},\qquad w_\beta=\frac{C_0-C_\alpha}{C_\beta-C_\alpha}.\end{equation}
Ambas fracciones suman uno y permanecen entre cero y uno dentro de la región bifásica. No son automáticamente fracciones volumétricas: convertirlas exige densidades. Fuera de esa región no se aplica la misma construcción.

Un enfriamiento rápido puede impedir que se alcance el equilibrio. Por eso un diagrama por sí solo no da dureza, tamaño de precipitado o fracción transformada tras un tiempo concreto. La cinética requiere información adicional; las curvas tiempo--temperatura--transformación y las de enfriamiento continuo responden a preguntas diferentes. No se emplean como calendarios intercambiables.
\ilustra{diagrama-fases}{Diagrama binario isomorfo esquemático con región líquido más sólido. Una línea de enlace en la región bifásica permite aplicar la regla de la palanca. No representa un sistema real ni permite definir un tratamiento industrial.}
\section{Precipitación: nucleación, crecimiento y sobreenvejecimiento}
Una secuencia típica de endurecimiento por precipitación parte de una solución sólida sobresaturada que evoluciona durante maduración. Precipitados finos y su relación con la matriz dificultan movimiento de dislocaciones. Al crecer y cambiar coherencia y separación puede variar el mecanismo de interacción y caer la resistencia tras un máximo. La explicación exige identificar la aleación: no todas las secuencias ni fases son iguales.

Una mayor dureza no demuestra automáticamente mejor tenacidad, corrosión bajo tensión o estabilidad térmica. El tiempo de servicio también puede actuar como exposición térmica y modificar el estado obtenido en fábrica. Si una reparación aplica calor local, el efecto puede extenderse más allá de la zona visible; no se restituye una condición solo por observar el mismo color de superficie.
\section{Control del proceso y de sus resultados}
Un proceso reproducible define identificación de pieza, estado inicial, equipos, medio, temperatura, permanencia, transferencia y enfriamiento, además de inspecciones posteriores. Las tolerancias y verificaciones proceden de una especificación aplicable; no se inventan a partir de un esquema docente. La evidencia combina registro del proceso y resultados pertinentes: dimensiones, dureza, microestructura u otros ensayos según la exigencia.
\repaso{Si se aumenta mucho el tiempo de maduración, ¿la resistencia debe seguir creciendo?}{No. Puede alcanzarse un máximo y después sobreenvejecimiento por evolución de precipitados. La respuesta depende de aleación, temperatura y estado previo; la dureza medida debe relacionarse con el requisito específico.}
